\documentclass[11pt,a4paper]{article}
\usepackage[margin=1in]{geometry}
\usepackage[T1]{fontenc}
\usepackage{lmodern}
\usepackage{setspace}
\setlength{\parindent}{0pt}
\setlength{\parskip}{6pt}
\linespread{1.05}

\begin{document}

\begin{center}
{\LARGE \textbf{EC7207: High Performance Computing}}\\[6pt]
{\large \textbf{Group Number: 36}}\\[6pt]
{\large \textbf{Parallel Protein Secondary Structure Prediction}}\\[8pt]
{\large Parallel Protein Secondary Structure Prediction on CB513 using Serial, OpenMP, POSIX Threads, MPI, and Hybrid (MPI+OpenMP)}
\end{center}

\textbf{Group Members}
\begin{itemize}
    \item EG/2021/4684 MUNSIF M.F.A.
    \item EG/2021/4685 MUTHUKUMARI H.M.S.
    \item EG/2021/4686 NADHA M.J.
\end{itemize}

\textbf{Project Summary}

This project addresses protein secondary structure prediction as a residue-level 3-class classification task (H/E/C) on the CB513 benchmark. We will build one common training pipeline and implement it across Serial, OpenMP, POSIX Threads, MPI, and Hybrid (MPI+OpenMP) models to evaluate performance and scalability on available CPU resources.

The dataset contains 513 protein files and 84,119 residues. DSSP labels from the raw CB513 files will be mapped to 3-state labels using a defined rule (H/G/I to H, B/E to E, others to C). Samples will be generated using a sliding window of length 13 with one-hot encoding (20 x 13 = 260 features), ensuring a fair, shared data pipeline for all implementations.

\textbf{Problem Statement}

Serial neural-network training on CB513 is too slow for repeated experiments and parameter sweeps. The key challenge is to parallelize training across shared-memory, distributed-memory, and hybrid settings while keeping prediction behavior consistent with the serial baseline. The project must therefore balance correctness (accuracy parity) and performance (timing/speedup) under real hardware constraints.

\textbf{Objectives}
\begin{enumerate}
    \item Build a correct serial baseline with reproducible preprocessing, training, and Q3 evaluation on CB513.
    \item Implement HPC variants with the same model/data settings: OpenMP, POSIX Threads, MPI, and Hybrid (MPI+OpenMP).
    \item Verify correctness by comparing each parallel variant against the serial baseline using Q3 accuracy.
    \item Measure runtime, speedup, and scaling by varying threads/ranks/other parameters under fixed hyperparameters.
    \item Prepare the final analysis report with a parallelization diagram and implementation description, accuracy comparison against serial, and timing comparison plots for serial vs parallel versions.
\end{enumerate}

\end{document}
